\section{Functor categories}\label{sec:functor_categories}

\begin{definition}\label{def:natural_transformation}\mcite[def. 1.3.1]{Leinster2014BasicCategoryTheory}
  Let \( F \) and \( G \) be parallel \hyperref[def:functor]{functors} from the category \( \cat{C} \) to \( \cat{D} \).

  A \term{natural transformation} \( \alpha \) from \( F \) to \( G \) is an \hyperref[def:indexed_family]{indexed family} of
  \begin{equation}\label{eq:def:natural_transformation/family}
    \seq{ \alpha_A: F(A) \to G(A) }_{A \in \cat{C}}
  \end{equation}
  of morphisms in \( \cat{D} \) such that, for every morphism \( f: A \to B \) in \( \cat{C} \), the following \hyperref[def:categorical_diagram]{diagram commutes}:
  \begin{equation}\label{eq:def:natural_transformation/diagram}
    \begin{aligned}
      \includegraphics[page=1]{output/def__natural_transformation}
    \end{aligned}
  \end{equation}

  The morphisms \( \alpha_A \) are called the components of \( \alpha \). We denote natural transformations by \( \alpha: F \Rightarrow G \) and, when used in diagrams, by
  \begin{equation}\label{eq:def:natural_transformation/notation}
    \begin{aligned}
      \includegraphics[page=2]{output/def__natural_transformation}
    \end{aligned}
  \end{equation}
\end{definition}

\begin{example}\label{ex:directed_multigraphs_as_functors}\mcite[example 1.3.46]{Perrone2024StartingCategoryTheory}
  In \cref{def:directed_multigraph}, we have defined a directed multigraph as a set \( V \) of vertices, a set \( A \) of arcs and two functions --- the initial \( i: A \to V \) and terminal \( t: A \to V \) vertices of an arc.

  Now consider the following \hyperref[def:categorical_diagram]{index category} \( \cat{I}: \)
  \begin{equation}\label{eq:ex:directed_multigraphs_as_functors/index/dots}
    \begin{aligned}
      \includegraphics[page=1]{output/ex__directed_multigraphs_as_functors}
    \end{aligned}
  \end{equation}

  For the sake of readability, we will give the following explicit labels in this category:
  \begin{equation}\label{eq:ex:directed_multigraphs_as_functors/index/annotated}
    \begin{aligned}
      \includegraphics[page=2]{output/ex__directed_multigraphs_as_functors}
    \end{aligned}
  \end{equation}

  A directed multigraph can then be defined as a functor \( G: \cat{I} \to \cat{Set} \) to the category \hyperref[def:category_of_small_sets]{\( \cat{Set} \)} of \hyperref[def:small_set]{small} sets.

  A \hyperref[def:natural_transformation]{natural transformation} from the directed multigraph \( G: \cat{I} \to \cat{Set} \) to \( H: \cat{I} \to \cat{Set} \) is then a pair of functions \( f_V: G(V) \to H(V) \) and \( f_A: G(A) \to H(A) \) such that the following \hyperref[def:categorical_diagram]{diagrams commute}:
  \begin{equation}\label{eq:ex:directed_multigraphs_as_functors/index/diagram}
    \begin{aligned}
      \includegraphics[page=3]{output/ex__directed_multigraphs_as_functors}
      \quad\quad\quad\quad
      \includegraphics[page=4]{output/ex__directed_multigraphs_as_functors}
    \end{aligned}
  \end{equation}

  See \cref{ex:isomorphism_of_directed_multigraph_categories} for how these functors related to directed multigraphs as defined in \cref{def:directed_multigraph}.
\end{example}

\begin{remark}\label{rem:natural_transformations_into_set}
  Let \( \cat{C} \) be an arbitrary \hyperref[def:small_set]{small} category. A \hyperref[def:natural_transformation]{natural transformation} \( \alpha \) from \( F: \cat{C} \to \cat{Set} \) to \( G: \cat{C} \to \cat{Set} \) is then a family of functions
  \begin{equation*}
    \seq{ \alpha_A: F(A) \to G(A) }_{A \in \cat{C}}.
  \end{equation*}

  Suppose that for every two objects \( A \) and \( B \) in \( \cat{C} \), the functions \( \alpha_A \) and \( \alpha_B \) agree on \( F(A) \cap F(B) \). This is automatically satisfied in \( F(A) \) and \( F(B) \) are disjoint whenever \( A \neq B \).

  We can then take the set-theoretic union of \( \alpha \) to obtain the function
  \begin{equation*}
    \bigcup_{A \in \cat{C}} \alpha_A: \bigcup\set{ F(A) \given A \in \cat{C} } \to \bigcup\set{ G(A) \given A \in \cat{C} }.
  \end{equation*}

  Both the domain and codomain are sets as a consequence of \ref{def:grothendieck_universe/union}, therefore the function is well-defined in the \hyperref[def:smallest_grothendieck_universe]{smallest Grothendieck universe} \( \mscrU_0 \). Denote it on \( \Alpha \) for brevity.

  An advantage of this is that we can define a natural transformation to be a function on a general enough set and then prove that its restrictions satisfy \eqref{eq:def:natural_transformation/diagram}.

  For example, consider the power set functor \( \pow: \cat{Set} \to \cat{Set} \) discussed in \cref{ex:def:functor}. The \hyperref[def:set_valued_map/identity]{identity function} \( \id_{\mscrU_0} \) is then a natural transformation from the identity functor \( \id_{\cat{Set}} \) to \( \pow \).

  Another natural transformation between the same functors is the singleton set operation \( \Sigma \) on sets defined as \( A \mapsto \set{ A } \). Note that, in this context, \( \Sigma \) operates not on the sets \( \id_{\cat{Set}}(A) \) and \( \pow(A) \), but on their members. The diagram \eqref{eq:def:natural_transformation/diagram} becomes
  \begin{equation}\label{eq:rem:natural_transformations_into_set}
    \begin{aligned}
      \includegraphics[page=1]{output/rem__natural_transformations_into_set}
    \end{aligned}
  \end{equation}

  This diagram commutes because, for every function \( f: A \to B \) and every \( x \in A \), we have
  \begin{equation*}
    f[\set{ x }] = \set{ f(x) }.
  \end{equation*}
\end{remark}

\begin{definition}\label{def:functor_category}
  Let \( \cat{C} \) and \( \cat{D} \) be arbitrary \hyperref[def:category]{categories}. The \term{functor category} \( [\cat{C}, \cat{D}] \), also denoted as \( \cat{D}^{\cat{C}} \), is defined as follows:

  \begin{itemize}
    \item The \hyperref[def:category/objects]{set of objects} \( \obj([\cat{C}, \cat{D}]) \) is the set of all functors from \( \cat{C} \) to \( \cat{D} \).

    \item The \hyperref[def:category/morphisms]{set of morphisms} \( [\cat{C}, \cat{D}](F, G) \) from \( F \) to \( G \) is the set of all \hyperref[def:natural_transformation]{natural transformations} from \( F \) to \( G \).

    \item The \hyperref[def:category/composition]{composition of the morphisms} \( \alpha: F \Rightarrow G \) and \( \beta: G \Rightarrow H \) is the natural transformation \( \beta \bincirc \alpha: F \Rightarrow H \) defined in terms of componentwise morphism composition, i.e.
    \begin{equation}\label{eq:def:functor_category/composition}
      (\beta \bincirc \alpha)_A \coloneqq \beta_A \bincirc \alpha_A.
    \end{equation}

    \item The \hyperref[def:category/identity]{identity morphism} on the functor \( F: \cat{C} \to \cat{D} \) is the \term{identity natural transformation} \( \id_F: F \Rightarrow F \) with components
    \begin{equation}\label{eq:def:functor_category/identity}
      (\id_F)_A \coloneqq \underbrace{\id_{F(A)}}_{F(\id_A)}
    \end{equation}
  \end{itemize}
\end{definition}
\begin{defproof}
  Just to verify that the composition \( \beta \bincirc \alpha \) defined in \eqref{eq:def:functor_category/composition} is indeed a natural transformation from \( F \) to \( H \), note that the following diagram trivially commutes:
  \begin{equation}\label{def:functor_category/composition}
    \begin{aligned}
      \includegraphics[page=1]{output/def__functor_category}
    \end{aligned}
  \end{equation}

  Now, to see that \( [\cat{C}, \cat{D}] \) is indeed a category, we verify the conditions \ref{def:category/C1} and \ref{def:category/C2}, which are in turn inherited from the same conditions on the categories \( \cat{C} \) and \( \cat{D} \).

  \SubProofOf{def:category/C1} For every two functors \( F, G: \cat{C} \to \cat{D} \) and natural transformation \( \alpha: F \Rightarrow G \), for every object \( A \in \cat{C} \) we have
  \begin{equation*}
    \id_{G(A)} \bincirc \alpha_A
    \reloset{\eqref{def:category/C1}} =
    \alpha_A
    \reloset{\eqref{def:category/C1}} =
    \alpha_A \bincirc \id_{F(A)}
  \end{equation*}

  Therefore,
  \begin{equation*}
    \id_G \bincirc \alpha = \alpha = \alpha \bincirc \id_F
  \end{equation*}
  and, after generalizing, we obtain that \ref{def:category/C1} holds in \( [\cat{C}, \cat{D}] \).

  \SubProofOf{def:category/C2} For any quadruple \( F \), \( G \), \( H \) and \( T \) of functors from \( \cat{C} \) to \( \cat{D} \) and every combination of natural transformations \( \alpha: F \Rightarrow G \), \( \beta: G \Rightarrow H \) and \( \gamma: H \Rightarrow T \), for every object \( A \in \cat{C} \) we have
  \begin{equation*}
    (\gamma_A \bincirc \beta_A) \bincirc \alpha_A
    \reloset{\eqref{def:category/C2}} =
    \gamma_A \bincirc (\beta_A \bincirc \alpha_A).
  \end{equation*}

  Therefore, after generalizing, we obtain that \ref{def:category/C2} holds in \( [\cat{C}, \cat{D}] \).
\end{defproof}

\begin{example}\label{ex:isomorphism_of_directed_multigraph_categories}
  In \cref{ex:directed_multigraphs_as_functors}, we defined \hyperref[def:directed_multigraph]{directed multigraphs} as functors from a certain index category \( \cat{I} \) to \( \cat{Set} \) (for a \hyperref[def:small_category]{fixed Grothendieck universe} \( \mscrU \)).

  There is then an obvious correspondence between directed multigraphs as objects of the category \( \cat{DM} \) defined in \cref{def:directed_multigraph/category}, and directed multigraphs as objects in the \hyperref[def:functor_category]{functor category} \( [\cat{I}, \cat{Set}] \), defined in \cref{ex:directed_multigraphs_as_functors}. Indeed, given any functor \( G: \cat{I} \to \cat{Set} \), the quadruple
  \begin{equation*}
    \parens[\big]{ G(V), G(A), G(h), G(t) }
  \end{equation*}
  is a directed multigraph in the sense of \cref{def:directed_multigraph}.

  No object in \( \cat{DM} \) is formally equal to any object in \( [\cat{I}, \cat{Set}] \) in the sense of \hyperref[def:zfc]{\logic{ZFC}}. They are, however, equivalent, as shown above, and this can be formalized by stating that the two categories are isomorphic, in the sense of \cref{def:morphism_invertibility/isomorphism}, as objects of the category \( \cat[\mscrV]{Cat} \), where \( \mscrV \) is a Grothendieck universe that strictly contains \( \mscrU \). We have already defined this isomorphism explicitly.

  This is an example of \term{isomorphism of categories}. In practice, if two categories are not so obviously identical, we are usually better served by \term{equivalences of categories} defined in \cref{def:category_equivalence}.
\end{example}

\begin{definition}\label{def:opposite_natural_transformation}\mimprovised
  The \term{opposite} natural transformation of \( \alpha: F \Rightarrow G \), where \( F \) and \( G \) are functors from \( \cat{C} \) to \( \cat{D} \), is the natural transformation \( \alpha^\oppos: G^\oppos \Rightarrow F^\oppos \), in which we take the opposite of each component in \( \alpha \).

  Dual natural transformation arise naturally in \cref{thm:opposite_of_functor_category}.
\end{definition}
\begin{defproof}
  The naturality diagram \eqref{eq:def:natural_transformation/diagram} commutes for \( \alpha^\oppos \) because all morphisms are simply reversed.
\end{defproof}

\begin{proposition}\label{thm:opposite_of_functor_category}
  For the \hyperref[def:opposite_category]{opposite} of the \hyperref[def:functor_category]{functor category} \( [\cat{C}, \cat{D}] \) we have
  \begin{equation*}
    [\cat{C}, \cat{D}]^\oppos = [\cat{C}^\oppos, \cat{D}^\oppos].
  \end{equation*}

  This is part of the duality principles listed in \fullref{thm:category_duality}.
\end{proposition}
\begin{proof}
  In \cref{def:opposite_functor}, we have defined the opposite functor \( F^\oppos: \cat{C}^\oppos \to \cat{D}^\oppos \) of \( F: \cat{C} \to \cat{D} \) in a way that allows us to regard it as an object of \( [\cat{C}^\oppos, \cat{D}^\oppos] \).

  In \cref{def:opposite_natural_transformation}, we have defined the opposite natural transformation \( \alpha^\oppos: G^\oppos \to F^\oppos \) of \( \alpha: F \Rightarrow G \) in a way that allows us to regard it as a morphism of \( [\cat{C}^\oppos, \cat{D}^\oppos] \).

  Furthermore, \( \alpha^\oppos: G^\oppos \to F^\oppos \) reverses the direction of its morphisms, and hence it is the dual to \( \alpha \) in the category \( [\cat{C}, \cat{D}]^\oppos \).
\end{proof}

\begin{definition}\label{def:diagonal_functor}\mcite[exerc. 3.1.1]{Leinster2014BasicCategoryTheory}
  Given an \term{index category} \( \cat{I} \) and an arbitrary category \( \cat{C} \), for any object \( A \) in \( \cat{C} \), we can define the \term{constant functor}
  \begin{equation*}
    \begin{aligned}
      &\Delta_A^{\cat{I}}: \cat{I} \to \cat{C}, \\
      &\Delta_A^{\cat{I}}(X) \coloneqq X, \\
      &\Delta_A^{\cat{I}}(g: X \to Y) \coloneqq \id_A.
    \end{aligned}
  \end{equation*}

  Given two objects \( A \) and \( B \) in \( \cat{C} \), a natural transformation \( \alpha: \Delta_A^{\cat{I}} \Rightarrow \Delta_B^{\cat{I}} \) is an \hyperref[def:indexed_family]{indexed family} that gives the same morphism for every object of the index category \( \cat{I} \).

  Indeed, the diagram \eqref{eq:def:natural_transformation/diagram} in this case becomes
  \begin{equation}\label{eq:def:diagonal_functor/nat}
    \begin{aligned}
      \includegraphics[page=1]{output/def__diagonal_functor}
    \end{aligned}
  \end{equation}

  This diagram implies that \( \alpha_A = \alpha_B \) for any two objects \( A \) and \( B \) in \( \cat{I} \). Therefore, all components of \( \alpha \) are equal to some morphism in \( \cat{C}(A, B) \).

  We can now define the \( \cat{I} \)-shaped \term{diagonal functor} on \( \cat{C} \)
  \begin{equation*}
    \begin{aligned}
      &\Delta^{\cat{I}}: \cat{C} \to [\cat{I}, \cat{C}], \\
      &\Delta^{\cat{I}}(A) \coloneqq \Delta_A^{\cat{I}}, \\
      &\Delta^{\cat{I}}(f: A \to B) \coloneqq \seq{ f: A \to B }_{k \in \cat{I}}.
    \end{aligned}
  \end{equation*}

  It is called a diagonal functor because, if \( \cat{I} \) is a discrete category of small two objects, then \( \Delta_{\cat{I}} \) gives the diagonal of the \hyperref[def:product_category]{product category} \( \cat{C}^2 \) by providing, for each object \( A \) of \( \cat{C} \), the ordered pair \( (A, A) \) (and similarly for morphisms).
\end{definition}

\begin{proposition}\label{thm:natural_isomorphism}\mcite{math3ma:natural_transformations}
  Let \( F \) and \( G \) be parallel \hyperref[def:functor]{functors} from the category \( \cat{C} \) to \( \cat{D} \). The family \eqref{eq:def:natural_transformation/family} is an isomorphism in the corresponding \hyperref[def:functor_category]{functor category} \( [\cat{C}, \cat{D}] \) if and only if all of its components are isomorphisms and, for any morphism \( f: A \to B \) in \( \cat{C} \), the following diagram commutes:
  \begin{equation}\label{eq:thm:natural_isomorphism/diagram}
    \begin{aligned}
      \includegraphics[page=1]{output/thm__natural_isomorphism}
    \end{aligned}
  \end{equation}

  We say that \( \alpha \) is a \term{natural isomorphism}.
\end{proposition}
\begin{proof}
  If all components of \( \alpha \) are isomorphisms, the condition
  \begin{equation*}
    \alpha_B \bincirc F(f) = G(f) \bincirc \alpha_A
  \end{equation*}
  is equivalent to
  \begin{equation*}
    F(f) = \alpha_B^{-1} \bincirc G(f) \bincirc \alpha_A.
  \end{equation*}

  We must now show that, if \( \alpha \) is an isomorphism in \( [\cat{C}, \cat{D}] \), all of its components are isomorphisms.

  If \( \alpha: F \Rightarrow G \) is an isomorphism in \( [\cat{C}, \cat{D}] \). Then there exists some natural transformation \( \beta: G \Rightarrow F \) such that
  \begin{equation*}
    \beta \bincirc \alpha = \id_F \quad\T{and}\quad \alpha \bincirc \beta = \id_G.
  \end{equation*}

  For every object \( A \) in \( \cat{C} \), the morphism \( \alpha_A: F(A) \to G(A) \) composed with \( \beta_A: G(A) \to F(A) \) is
  \begin{equation*}
    \beta_A \bincirc \alpha_A = \id_{F(A)}.
  \end{equation*}

  Therefore, \( \alpha_A \) is left invertible. Analogously,
  \begin{equation*}
    \alpha_A \bincirc \beta_A = \id_{F(A)}
  \end{equation*}
  and hence \( \alpha_A \) is right invertible.

  Therefore, for every object \( A \) in \( \cat{C} \), the morphism \( \alpha_A \) is fully invertible, i.e. an isomorphism.
\end{proof}

\begin{definition}\label{def:function_currying}\mcite[259]{TroelstraSchwichtenberg2000BasicProofTheory}
  Given the \hyperref[con:function_arguments]{binary} function \( f: A \times B \to C \), we can define another function \( g: A \to \fun(B, C) \) as
  \begin{equation*}
    g(x) \coloneqq (y \mapsto f(x, y)).
  \end{equation*}

  We call this process \term{currying} after Haskell Curry. Obtaining \( f \) from \( g \) can be called \term{uncurrying}.
\end{definition}
\begin{comments}
  \item Currying is useful if we have somehow fixed a value \( x_0 \in A \), in which case we can \enquote{get rid} of one argument by introducing some shortcut for the function \( g(x_0): B \to C \) for the sake of reducing notational clutter. See \cref{def:differentiability/first_variation} and our proof of \cref{thm:countably_infinite_union_of_countably_infinite_sets} for example of how this is useful in the wild.
\end{comments}

\begin{proposition}\label{thm:currying_is_natural_isomorphism}
  \hyperref[def:function_currying]{Function currying} is a natural isomorphism between the functors
  \begin{align*}
    &\cat{Set}(A \times B, C)
    &\cat{Set}(A, \cat{Set}(B, C))
  \end{align*}

  More concretely, consider the following functors, which are loosely based on \cref{thm:hom_functor}:
  \begin{equation*}
    \begin{aligned}
      &V: \cat{Set}^3 \to \cat{Set} \\
      &V(A, B, C) \coloneqq \cat{Set}(A \times B, C) \\
      \Big[& V(f: X \to A, g: Y \to B, h: C \to Z) \Big](s: A \times B \to C) \coloneqq \smash{ \overbrace{ (x, y) \mapsto h\parens[\Bigg]{ s\parens[\big]{ \underbrace{ f(x), g(y) }_{A \times B} } } }^{X \times Y \to Z} } \\
    \end{aligned}
  \end{equation*}
  and
  \begin{equation*}
    \begin{aligned}
      &W: \cat{Set}^3 \to \cat{Set} \\
      &W(A, B, C) \coloneqq \cat{Set}(A, \cat{Set}(B, C)) \\
      \Big[& W(f: X \to A, g: Y \to B, h: C \to Z) \Big](t: A \to \cat{Set}(B, C)) \coloneqq \smash{ \underbrace{ x \mapsto \overbrace{y \mapsto h\parens[\Bigg]{ \overbrace{t(f(x))}^{B \to C}\parens[\big]{ g(y) } } }^{Y \to Z} }_{X \to \cat{Set}(Y, Z)} } \\
    \end{aligned}
  \end{equation*}

  Then the family of functions
  \begin{equation*}
    \begin{aligned}
      &\alpha: V \Rightarrow W \\
      &\alpha_{A,B,C}(s: A \times B \to C) \coloneqq a \mapsto b \mapsto s(a, b)
    \end{aligned}
  \end{equation*}
  is a \hyperref[thm:natural_isomorphism]{natural isomorphism}.
\end{proposition}
\begin{proof}
  The function \( \varphi \) is clearly invertible. Fix a triple of functions \( f: X \to A \), \( g: Y \to B \) and \( h: C \to Z \). For every \( s: A \times B \to C \) we have
  \begin{align*}
    [W(f, g, h)](\alpha_{A,B,C}(s))
    &=
    x \mapsto y \mapsto \parens[\big]{ h\parens[\big]{ [a \mapsto b \mapsto s(a, b)](f(x))(g(y)) } }
    = \\ &=
    x \mapsto y \mapsto h\parens[\big]{ s(f(x), g(y)) }
    = \\ &=
    \alpha_{X,Y,Z}(V(f, g, h)),
  \end{align*}
  which proves that the following diagram commutes:
  \begin{equation}\label{eq:thm:currying_is_natural_isomorphism/diagram}
    \begin{aligned}
      \includegraphics[page=1]{output/thm__currying_is_natural_isomorphism}
    \end{aligned}
  \end{equation}
\end{proof}
